% Główny plik pracy magisterskiej (modułowy).
\documentclass{ppfcmthesis}

% ------------------- PAKIETY OGÓLNE -------------------
\usepackage{amsmath,amssymb}
\usepackage{microtype} % lepsza typografia (dzielenie, kerning)
\usepackage{setspace} % dla \setstretch w streszczeniach
\usepackage{siunitx}
\usepackage{listings}
\usepackage{graphicx}
\usepackage{subcaption}
\usepackage{xcolor}
\usepackage{booktabs}
\usepackage{multirow}
\usepackage{tabularx}
\usepackage{longtable}
\usepackage{float}
\usepackage{enumitem}
\usepackage{tikz}
\usetikzlibrary{arrows.meta,positioning,fit,shapes,calc}
% Algorytmy: unikamy pakietu 'algorithm' (ostrzeżenie UTF-8 w Tectonic)
\usepackage{algpseudocode} % zapewnia środowisko algorithmic
\usepackage{float}
\floatstyle{ruled}
\newfloat{algorithm}{tbp}{loa}
\floatname{algorithm}{Algorytm}
\usepackage{csquotes}

% Bibliografia (biblatex + biber)
\usepackage[
  backend=bibtex,
  style=ieee,
  doi=false,
  isbn=false,
  url=true,
  giveninits=true,
  sorting=nyt
]{biblatex}
\addbibresource{bibliografia.bib}

% Polskie nazwy spisów
\addto\captionspolish{%
  \renewcommand{\contentsname}{Spis treści}% (już jest, dla pewności)
  \renewcommand{\listfigurename}{Spis rysunków}%
  \renewcommand{\listtablename}{Spis tabel}%
}
% Nazwa listy listingów
\renewcommand{\lstlistlistingname}{Spis listingów}

% ------------------- LISTINGS -------------------
\lstdefinestyle{code}{
  basicstyle=\ttfamily\small,
  numbers=left,
  numberstyle=\tiny,
  numbersep=6pt,
  frame=single,
  framerule=0.4pt,
  breaklines=true,
  tabsize=2,
  showstringspaces=false,
  keywordstyle=\color{blue!70!black}\bfseries,
  commentstyle=\color{green!40!black}\itshape,
  stringstyle=\color{orange!60!black},
  captionpos=b
}
\lstset{style=code,numberbychapter=false}
% Prosta globalna numeracja listingów
\AtBeginDocument{\renewcommand{\thelstlisting}{\arabic{lstlisting}}}

% ------------------- NUMERACJA RÓWNAŃ / ALGORYTMÓW -------------------
\numberwithin{equation}{chapter}
\numberwithin{figure}{chapter}
\numberwithin{table}{chapter}
\numberwithin{algorithm}{chapter}
% Głębokość numerowania sekcji: rozdział (1), sekcja (1.1), podsekcja (1.1.1)
\setcounter{secnumdepth}{2}

% ------------------- MAKRA BADAWCZE / SKRÓTY -------------------
\newcommand{\cel}[1]{\textbf{Cel:} #1}
\newcommand{\pytanie}[1]{\textbf{Pytanie badawcze:} #1}
\newcommand{\hipoteza}[1]{\textbf{Hipoteza:} #1}
\newcommand{\abbr}[2]{\textbf{#1} – #2\\}

% ------------------- METADANE -------------------
\title{Analiza ruchu sieciowego aplikacji mobilnych na platformie Android}
\authortitle{Imię Nazwisko, nr albumu 123456}
\author{Imię Nazwisko\album{123456}}
\ppsupervisor{dr inż. Imię Nazwisko}
\ppyear{2025}

% (Opcja: identyczne logo na kolofonie)
% \renewcommand{\ppColophonLogoFile}{wit/sygnet-pp}
% \renewcommand{\ppColophonLogoWidth}{0.5\textwidth}

% (Opcja Times – jeśli zechcesz później)
% \usepackage{newtxtext}
% \usepackage{newtxmath}

\begin{document}
\maketitle

% ======= STRESZCZENIA =======
% 00-streszczenia.tex — oba streszczenia wyśrodkowane, jedno pod drugim
\clearpage
\thispagestyle{plain} % albo 'empty', jeśli chcesz bez numeru strony
\vspace*{\fill}

\begin{center}

  % ====== PL ======
  {\Large\bfseries Streszczenie\par}
  \vspace{0.6\baselineskip}
  \begin{minipage}{0.85\textwidth}
    \small
    \noindent
  Celem pracy jest zbudowanie powtarzalnego środowiska do analizy ruchu aplikacji Android pod kątem bezpieczeństwa transmisji i prywatności. \emph{Pipeline} łączy rejestrację na urządzeniu (\emph{PCAP/PCAPNG}), klasyfikację protokołów (\emph{TLS}, \emph{HTTP/1.1–3/QUIC}) oraz korelację ze statyczną analizą (Exodus). Badania pokazują dominację \emph{HTTP/3/QUIC} w aplikacjach społecznościowych, konserwatywne użycie \emph{TLS/HTTPS} i niewiele domen w aplikacjach instytucjonalnych oraz typową telemetrię i zaplecze chmurowe w komercyjnych; odnotowano też ryzyko funkcjonalne (scenariusz Posnania). Ograniczenia: zmienność wersji i pinning/\emph{QUIC}; rekomendacja: automatyczna kategoryzacja hostów i monitoring domen trzecich.
    
    \medskip
    \noindent\textbf{Słowa kluczowe:} Android, analiza ruchu, TLS, HTTP/3/QUIC, prywatność, PCAP, telemetria.
  \end{minipage}

  % przerwa między blokami (mała)
  \vspace{10\baselineskip}

  % ====== EN ======
  {\Large\bfseries Abstract\par}
  \vspace{0.6\baselineskip}
  \begin{minipage}{0.85\textwidth}
    \small
    \noindent
  This thesis builds a reproducible framework for analyzing Android app traffic with a focus on transport security and privacy. The pipeline combines on-device capture (\emph{PCAP/PCAPNG}), protocol classification (\emph{TLS}, \emph{HTTP/1.1–3/QUIC}), and correlation with static analysis (Exodus). Results show \emph{HTTP/3/QUIC} dominance in social/media apps, conservative \emph{TLS/HTTPS} with few third-party domains in institutional apps, and common telemetry plus cloud backends in commercial apps; a functional risk was observed (Posnania scenario). Limitations: version variability and pinning/\emph{QUIC}; recommendation: automatic host categorization and continuous third-party monitoring.
    
    \medskip
    \noindent\textbf{Keywords:} Android, network traffic analysis, TLS, QUIC, privacy, mitmproxy.
  \end{minipage}

\end{center}

\vspace*{\fill}
\clearpage


\tableofcontents*
\cleardoublepage
\listoffigures
\cleardoublepage
\listoftables
\cleardoublepage
\listof{algorithm}{Spis algorytmów}
\cleardoublepage
\lstlistoflistings
\cleardoublepage
\chapter*{Wykaz skrótów}
\addcontentsline{toc}{chapter}{Wykaz skrótów}
\abbr{API}{Application Programming Interface}
\abbr{TLS}{Transport Layer Security}
\abbr{HTTP}{HyperText Transfer Protocol}
\abbr{QUIC}{Quick UDP Internet Connections}
\abbr{CDN}{Content Delivery Network}
\abbr{DNS}{Domain Name System}
\abbr{ADB}{Android Debug Bridge}


\mainmatter*

% ======= ROZDZIAŁY GŁÓWNE =======
\chapter{Wstęp}

\section{Motywacja}

Rynek aplikacji mobilnych stał się dominującym kanałem konsumpcji treści i usług cyfrowych. Dane o penetracji dostępu mobilnego (ok.\ 71\% globalnej populacji posiada subskrypcję mobilną do 2023 r.) \cite{CiscoAIRWhitepaper2023} oraz statystyki udziału ruchu WWW (ponad 50\% globalnych \emph{page views} pochodzi z urządzeń mobilnych) \cite{StatCounterPlatformShare2025} pokazują, że to właśnie środowisko mobilne jest dziś kluczowym miejscem, w którym kształtują się wzorce komunikacji sieciowej i materializują ryzyka dla prywatności.

W ekosystemie Androida współistnieje wiele bibliotek i zestawów SDK firm trzecich — od analityki i reklam, przez personalizację, po rozwiązania antyfraudowe. Każda z tych warstw wnosi własne punkty końcowe, protokoły i mechanizmy telemetrii, często działające poza świadomością użytkownika, a niekiedy również poza pełną kontrolą zespołów deweloperskich aplikacji. W efekcie rośnie złożoność oraz nieprzejrzystość przepływów danych.

Konsekwencją tego stanu rzeczy są konkretne wyzwania badawcze:
\begin{enumerate}[label=\alph*)]
  \item uzyskanie przejrzystości, \emph{jakie} protokoły, domeny i dostawcy są realnie wykorzystywani przez aplikacje,
  \item ocena ryzyka ekspozycji metadanych (np.\ hosty, czasy, rozmiary transferów) mimo powszechnego szyfrowania treści,
  \item porównywalność zachowań komunikacyjnych między kategoriami aplikacji (komunikatory, multimedia, finanse, gry),
  \item budowa powtarzalnych narzędzi i procedur, które nadążą za ewolucją protokołów (HTTP/2, HTTP/3/QUIC) oraz mechanizmami obronnymi (np.\ \emph{certificate pinning}).
\end{enumerate}

Wobec powyższego potrzebne jest uporządkowane, możliwie obiektywne podejście do \emph{systematycznej} obserwacji i oceny ruchu sieciowego aplikacji: środowisko pomiarowe, które pozwala porównywać różne aplikacje i kategorie, zdefiniowany zestaw metryk (zarówno ilościowych, jak i jakościowych) oraz protokoły postępowania uwzględniające ograniczenia współczesnych technologii transportowych i zabezpieczeń. Tak rozumiana ramka badawcza stanowi bezpośrednią motywację dla niniejszej pracy.


Drugi wektor motywacyjny stanowi szybka adopcja nowych protokołów transportowych i warstw aplikacyjnych. QUIC (standaryzowany jako HTTP/3 \cite{RFC9114}) przenosi część logiki transportu do przestrzeni użytkownika (UDP + TLS 1.3 \cite{RFC8446}), co zmienia charakterystykę obserwowalnych cech przepływów: klasyczne estymatory opóźnień i retransmisji z poziomu TCP stają się mniej bezpośrednie, a analiza pakietów jest utrudniona przez szyfrowanie większej liczby pól. W konsekwencji wcześniejsze proporcje protokołów publikowane w literaturze mogą być nieaktualne dla obecnych wersji popularnych aplikacji.

Kolejnym aspektem jest prywatność i otoczenie regulacyjne. RODO (GDPR), a także DMA/DSA w UE akcentują minimalizację i przejrzystość przetwarzania. Architektura „multi-endpoint” (API własne + analityka + reklamy + antyfraud + CDN/edge) prowadzi do gęstych, równoległych zestawów połączeń tuż po interakcji użytkownika. Nawet przy pełnym szyfrowaniu treści korelacja metadanych (częstotliwość, rozmiary odpowiedzi, domeny) może ujawniać wzorce użycia funkcji (pattern-of-life). Transparentność takich wzorców jest ograniczona zarówno dla użytkownika, jak i audytorów. Narzędzia systemowe dają uogólniony obraz, a techniki inwazyjne (hooking, dynamic instrumentation) bywają blokowane przez mechanizmy anty‑tampering.

Istnieje także motywacja praktyczna stworzenia lekkiego, powtarzalnego środowiska badawczego balansującego między:
\begin{itemize}
  \item \textbf{Zasięgiem obserwacji} (kompletność metadanych HTTP[S]),
  \item \textbf{Inwazyjnością} (brak konieczności rootowania, minimalne modyfikacje),
  \item \textbf{Zgodnością prawną i etyczną} (anonimizacja identyfikatorów, brak utrwalania treści wrażliwych).
\end{itemize}
Istniejące narzędzia (np. \emph{mitmproxy}) zapewniają bazowy przechwyt, ale wymagają rozszerzeń o: 
\begin{itemize}
  \item \textbf{Normalizacja} – ujednolicenie formatu danych.
  \item \textbf{Klasyfikacja} – kategoryzacja ruchu sieciowego.
  \item \textbf{Agregacja} – zbieranie i analiza metryk.
  \item \textbf{Automatyzacja} – skrypty do interakcji z aplikacjami.
  \item \textbf{Anonimizacja} – ochrona danych osobowych.
\end{itemize}

Z naukowego punktu widzenia identyfikowane luki obejmują:
\begin{description}
  \item[L1] Brak bieżących, empirycznych pomiarów dystrybucji HTTP/1.1 vs HTTP/2 vs QUIC w popularnych aplikacjach.
  \item[L2] Ograniczone porównania między kategoriami aplikacji a zachowaniami komunikacyjnymi.
  \item[L3] Niedostatek jawnych metryk ryzyka prywatności bazujących wyłącznie na metadanych.
  \item[L4] Fragmentaryczne dane o wpływie certificate pinning na pokrycie przechwytu.
\end{description}

Interesariusze:
\begin{itemize}
  \item \textbf{Użytkownik} – zyskuje pośrednio poprzez privacy-by-design.
  \item \textbf{Deweloper} – możliwość audytu nadmiarowych zależności.
  \item \textbf{Audytor bezpieczeństwa} – metryki priorytetyzacji ryzyka.
  \item \textbf{Społeczność naukowa} – powtarzalny pipeline do dalszych badań.
\end{itemize}

\section{Stan wiedzy}

\subsection{Analiza ruchu sieciowego}
Badania pokazują, że ruch generowany przez aplikacje mobilne niesie znaczące wyzwania dla prywatności – nawet konfiguracje TLS nie chronią całkowicie, bo analizę można przeprowadzać na metadanych lub kanałach bocznych. W przeglądzie literatury poświęconym analizie ruchu mobilnego autorzy klasyfikują cele, miejsca przechwytywania oraz metody wykorzystania danych sieciowych do profilowania użytkowników \cite{Conti2017TrafficAnalysis}. Inne prace demonstrują, jak sieć może zostać wykorzystana do identyfikacji konkretnej aplikacji, nawet przy szyfrowaniu, wykorzystując uczenie maszynowe na podstawie cech pakietów \cite{Taylor2017AppFingerprinting}. Przykłady wykorzystania narzędzi takich jak \textit{mitmproxy} i \textit{tcpdump} pokazują, jak można wychwycić i analizować zarówno HTTP jak i inne kanały sieciowe \cite{mitmproxy,tcpdump}.

\subsection{Inżynieria wsteczna i testy penetracyjne}
Inżynieria wsteczna aplikacji mobilnych, ułatwiająca badanie mechanizmów takich jak SSL pinning, jest szeroko stosowana. Wczesne prace wskazują na powszechne błędy konfiguracji SSL w aplikacjach Androidowych — ponad połowa analizowanych aplikacji korzystających z SSL miała błędne wdrożenia lub brak obsługi poprawnego pinningu \cite{Fahl2012WhyEve}. W celu przeprowadzania dynamicznych testów bezpieczeństwa wykorzystuje się narzędzia takie jak \textit{Frida}, \textit{Drozer} czy \textit{MobSF}, które wspierają zarówno analizę dynamiczną, jak i automatyzację \cite{MobSF,Frida}.

\subsection{Prywatność, uprawnienia i trackery}
Coraz więcej badań zwraca uwagę na skalę gromadzenia danych przez aplikacje mobilne. Systematyczne opracowania przeglądowe pokazują spektrum metod oceny prywatności (statycznych, dynamicznych i hybrydowych) oraz wskazują luki w zgodności z regulacjami \cite{DelAlamo2021PrivacyMapping}. Narzędzia takie jak \textit{Exodus Privacy} umożliwiają analizę APK pod kątem wbudowanych trackerów i żądanych uprawnień \cite{ExodusPrivacy}. Wyniki badań empirycznych pokazują, że aplikacje często żądają nadmiarowych uprawnień lub integrują zewnętrzne trackery reklamowe \cite{Razaghpanah2018AppsPrivacy}.

\subsection{Zużycie zasobów i kontekst behawioralny}
Nadmierne użycie zasobów (np. bateria, CPU) może wskazywać na obecność komponentów śledzących lub brak optymalizacji. Do pomiarów wykorzystuje się m.in. \textit{Battery Historian} oraz \textit{Android Profiler} \cite{BatteryHistorian,AndroidProfiler}. Chociaż ten aspekt jest rzadziej eksplorowany w literaturze naukowej, podkreśla się jego znaczenie w identyfikacji ukrytych mechanizmów aplikacji \cite{Pathak2012BatteryDrain}.

\subsection{Podsumowanie}
Przegląd literatury wykazuje, że istnieje bogata baza wiedzy w obszarach analizy ruchu sieciowego, inżynierii wstecznej, testów penetracyjnych oraz badania prywatności aplikacji mobilnych. Jednocześnie wciąż brakuje holistycznego podejścia, które łączyłoby te wszystkie elementy — co stanowi niszę, którą niniejsza praca ma ambicję wypełnić.


\section{Cel i zakres pracy}
\cel{Celem pracy jest opracowanie i ewaluacja środowiska badawczego umożliwiającego systematyczną analizę bezpieczeństwa i prywatności aplikacji mobilnych na system Android. Praca koncentruje się na przechwytywaniu, klasyfikacji oraz ocenie ruchu sieciowego, a także na badaniu mechanizmów uwierzytelniania, uprawnień i zużycia zasobów.}

Zakres pracy obejmuje:
\begin{itemize}
  \item \textbf{Środowisko przechwytywania} – konfiguracja narzędzi takich jak \textit{mitmproxy}, \textit{Wireshark}, \textit{tcpdump}, a także przygotowanie środowiska laboratoryjnego z wykorzystaniem emulatorów oraz zrootowanych urządzeń mobilnych.
  \item \textbf{Analiza ruchu sieciowego} – identyfikacja używanych protokołów (HTTP/1.1, HTTP/2, QUIC, TLS), badanie bezpieczeństwa transmisji danych (szyfrowanie, SSL pinning) oraz klasyfikacja przesyłanych informacji.
  \item \textbf{Reverse engineering i testy penetracyjne} – dekompilacja aplikacji z wykorzystaniem \textit{APKTool}, \textit{Jadx}, \textit{Ghidra} oraz testy przy użyciu narzędzi takich jak \textit{MobSF}, \textit{Frida}, \textit{Metasploit}.
  \item \textbf{Analiza uprawnień i prywatności} – badanie przydzielanych uprawnień, wykrywanie trackerów i komponentów reklamowych (\textit{Exodus Privacy}, \textit{App Inspector}) oraz ocena zarządzania danymi użytkowników.
  \item \textbf{Analiza zużycia zasobów} – pomiary obciążenia baterii i CPU aplikacji z użyciem \textit{Battery Historian} oraz \textit{Android Profiler}.
  \item \textbf{Eksperymenty} – testy na wybranych aplikacjach z różnych kategorii (m.in. komunikatory, aplikacje społecznościowe, bankowe, rządowe), analiza różnic w implementacji mechanizmów bezpieczeństwa.
  \item \textbf{Badanie wektorów ataków} – symulacja ataków typu \textit{Man-in-the-Middle}, analiza przechowywania danych lokalnych, scenariusze działania offline/online, testy odporności na automatyczne skrypty i boty.
  % \item \textbf{Analiza logów i ścieżki audytu} – sprawdzenie, jakie dane aplikacje rejestrują w logach oraz ocena ryzyka wycieku informacji.
\end{itemize}

\chapter{Podstawy teoretyczne}

\section{Architektura systemu Android}
Android jest otwartą, wielowarstwową platformą opartą na jądrze Linuksa. Stos oprogramowania obejmuje m.in.\ jądro (kernel) z mechanizmami bezpieczeństwa (np.\ SELinux), warstwę HAL (Hardware Abstraction Layer), biblioteki natywne, środowisko uruchomieniowe ART oraz framework Java/Kotlin z zestawem API, na którym budowane są aplikacje i usługi systemowe \cite{AOSPArchitecture,AndroidPlatformArchitecture}. W nowszych wersjach platformy pojawiły się dodatkowe izolatory, takie jak \emph{SDK sandbox} (separowany proces dla wybranych SDK), co jeszcze bardziej ogranicza wpływ komponentów zewnętrznych na proces aplikacji \cite{SdkSandbox}.

\section{Model bezpieczeństwa}
Fundamentem bezpieczeństwa Androida jest \emph{piaskownica} (application sandbox): każda aplikacja otrzymuje unikalny UID i działa w odseparowanym procesie z izolowanym przestrzennie systemem plików. Dodatkową ochronę zapewnia polityka MAC (SELinux), a także model uprawnień ograniczający dostęp do wrażliwych zasobów (lokalizacja, kamera, mikrofon) – w tym uprawnienia żądane w czasie działania (\emph{runtime permissions}) \cite{AndroidSandbox,AndroidPermissionsOverview,AndroidPermissionRequest}. Łańcuch zaufania domykają mechanizmy podpisywania aplikacji, Verified Boot/AVB oraz proces weryfikacji w ekosystemie sklepu (Play Protect, audyty polityk), co Google cyklicznie raportuje \cite{GooglePlaySafety2024}. 
W kontekście komunikacji sieciowej część aplikacji implementuje \emph{certificate pinning}, które utrudnia badaczowi pasywne przechwycenie treści HTTPS przez klasyczne proxy – poprawnie wdrożony pinning ma powodować odrzucenie połączenia z certyfikatem wystawionym przez zaufane, lecz nieprawidłowe (dla danej domeny) CA \cite{OWASPMASTGPinning}.

\section{Protokoły}
Przeglądane aplikacje korzystają z rodziny protokołów HTTP w różnych wersjach oraz z TLS do uwierzytelnienia i poufności. 
\begin{itemize}
  \item \textbf{HTTP/1.1} – klasyczny, tekstowy protokół nad TCP (przykładowe różnice semantyczne i transportowe względem nowszych wersji omawiają poniższe punkty HTTP/2/HTTP/3). 
  \item \textbf{HTTP/2} – binarne ramkowanie, kompresja nagłówków i \emph{multiplexing} wielu strumieni na jednym połączeniu TCP w celu redukcji opóźnień \cite{RFC9113}.
  \item \textbf{HTTP/3} – mapowanie semantyki HTTP na QUIC (UDP), eliminujące wpływ HoL na poziomie transportu i skracające zestawienie połączenia \cite{RFC9114,RFC9000}. 
  \item \textbf{TLS 1.3} – domyślny protokół kryptograficzny dla HTTPS, skracający \emph{handshake}, upraszczający zestaw szyfrów i wspierający 0-RTT (z zastrzeżeniami bezpieczeństwa) \cite{RFC8446}.
\end{itemize}

\section{Techniki analizy ruchu}
W analizach bezpieczeństwa mobilnego wykorzystuje się zwykle dwa komplementarne podejścia: 
\emph{(1) przechwytywanie i inspekcję ruchu} oraz \emph{(2) instrumentację/dynamiczną obserwację aplikacji}. 
Do przechwytywania pakietów i strumieni wykorzystywane są m.in.\ \textit{tcpdump} i \textit{Wireshark}, a do pośredniczenia w ruchu HTTP(S) – \textit{mitmproxy} (instalacja zaufanego CA na urządzeniu w celu deszyfracji ruchu) \cite{tcpdump,WiresharkDocs,MitmproxyDocs,OWASPMASTGIntercept}. 
W przypadku aplikacji z poprawnym \emph{certificate pinning} wymagane bywa podejście z instrumentacją (np.\ \textit{Frida}) lub środowisko z większą kontrolą (urządzenie z dostępem \textit{root}/emulator, własny trust store), przy czym nowoczesne techniki (np.\ niestandardowi \emph{trust managerzy}, natywne weryfikacje, integracja z \emph{Play Integrity}) istotnie podnoszą poprzeczkę dla badań dynamicznych \cite{Frida,OWASPMASTGPinning}. 
Wybór między emulatorem a urządzeniem fizycznym wpływa na wierność zachowania (emulator ułatwia kontrolę i automatyzację; urządzenie fizyczne lepiej oddaje ograniczenia sprzętowe i integracje OEM) \cite{OWASPMASTGIntercept}.

\section{Prywatność i tracking}
Problem prywatności w aplikacjach mobilnych obejmuje zarówno zakres i uzasadnienie uprawnień, jak i obecność komponentów stron trzecich (SDK analityczne/reklamowe). Do statycznego wykrywania trackerów i przeglądu uprawnień w APK wykorzystywane są m.in.\ \textit{Exodus Privacy} oraz narzędzia deweloperskie Androida \cite{ExodusPrivacy,AndroidPermissionsOverview}. Równolegle platforma rozwija rozwiązania mające ograniczyć śledzenie międzyaplikacyjne (np.\ \emph{Privacy Sandbox on Android}) \cite{PrivacySandboxAndroid}. W praktyce ocena prywatności powinna łączyć analizę statyczną i dynamiczną (ruch sieciowy, zachowanie w tle), a także kontrolować zużycie zasobów (bateria/CPU) – symptomy nadmiernej telemetrii mogą ujawniać się jako wzrost aktywności sieciowej w tle czy anomalia energetyczne.

\chapter{Metodyka badawcza}

\section{Kryteria doboru aplikacji}
Popularność (rankingi store), kategorie (komunikacja, multimedia, finanse), aktualność wersji.

\section{Środowisko testowe}
Sprzęt (model urządzenia), wersja Android, konfiguracja sieci (oddzielony SSID), kontrola tła (tryb samolotowy przed startem sesji).

\section{Procedura przechwytywania}
Instalacja certyfikatu CA dla mitmproxy, konfiguracja proxy systemowego, ograniczenia w wypadku pinningu.

\section{Automatyzacja interakcji}
Skrypty ADB / uiautomator generujące sekwencje działań, logowanie timestampów.

\section{Model danych i metryki}
Pola: timestamp, host, ścieżka, metoda, kod statusu, rozmiary (request/response), protokół (HTTP/2 / QUIC), kategoria domeny, czas odpowiedzi.

\section{Anonimizacja i etyka}
Hashowanie identyfikatorów, odrzucanie payloadów binarnych zawierających potencjalne dane osobowe, brak logowania treści formularzy.

\section{Walidacja}
Losowa weryfikacja spójności liczby pakietów vs liczby zarejestrowanych żądań; test powtarzalności (2 przebiegi).
\chapter{Architektura rozwiązania}

\section{Przegląd komponentów}
Opis modułów: przechwytywanie, ekstrakcja, normalizacja, klasyfikacja, agregacja, wizualizacja.

\section{Diagram wysokopoziomowy}
\begin{figure}[H]
  \centering
  % (Możesz wstawić diagram TikZ lub grafikę.)
  \caption{Architektura modułowa środowiska analitycznego (schemat).}
  \label{fig:architektura}
\end{figure}

\section{Przepływ danych (ETL)}
Kolejne etapy: capture -> parse -> enrich -> store -> stats.

\section{Struktura repozytorium}
Opis katalogów (np. scripts/, data/, analysis/).

\section{Model normalizacji}
Mapowanie surowych atrybutów mitmproxy na ujednolicony rekord JSON/CSV.

\section{Rozszerzalność}
Dodanie nowego klasyfikatora domen lub obsługa innego proxy.
\chapter{Implementacja}

\section{Moduł przechwytywania}
Konfiguracja mitmproxy, parametry uruchomienia.

\section{Hooki mitmproxy}
\begin{lstlisting}[language=Python,caption={Hook mitmproxy zapisujący metadane żądań},label={lst:hook}]
def response(flow):
    rec = {
        "proto": flow.request.scheme,
        "host": flow.request.host,
        "path": flow.request.path,
        "method": flow.request.method,
        "status": flow.response.status_code,
        "t_start": flow.request.timestamp_start,
        "t_end": flow.request.timestamp_end,
        "length": len(flow.response.raw_content or b"")
    }
    save_record(rec)
\end{lstlisting}

\section{Parser i normalizacja}
Krok po kroku: walidacja, cast typów, uzupełnianie braków.

\section{Klasyfikacja domen}
Listy znanych trackerów, heurystyki (substring, kategorie TLD, reverse DNS).

\section{Agregacje i raporty}
Obliczanie histogramów, percentyle czasów odpowiedzi.

\section{Testy}
Testy jednostkowe (poprawność mapowania), testy integracyjne (ciąg ETL).
\chapter{Eksperymenty i wyniki}

\section{Plan eksperymentalny}
Opis serii sesji, długość, liczba powtórzeń.

\section{Charakterystyka próbek}
Liczba aplikacji, łączna liczba żądań, udział kategorii.

\section{Wykorzystanie protokołów}
Tabela i wykres udziału HTTP/1.1 vs HTTP/2 vs QUIC.

\section{Analiza domen trackingowych}
Odsetek żądań do trackerów per kategoria.

\section{Czasy odpowiedzi i rozmiary}
Rozkłady, mediany, percentyle.

\section{Ograniczenia}
Próbkowanie krótkich okien użycia, brak tła push/pull w długim horyzoncie.
\chapter{Aspekty bezpieczeństwa i prywatności}

\section{Szyfrowanie i wersje TLS}
Wykryte wersje, renegocjacje (jeśli obserwowane).

\section{Certificate pinning}
Przypadki utraty widoczności ruchu (fail-open / fail-closed), wpływ na wyniki.

\section{Potencjalne wycieki metadanych}
Parametry w URL, identyfikatory urządzeń, tokeny.

\section{Domeny trackingowe}
Porównanie intensywności ruchu do domen analitycznych.

\section{Implikacje}
Ryzyka dla prywatności użytkownika i rekomendacje (np. ograniczanie uprawnień sieciowych, firewall lokalny).
\chapter{Dyskusja}

\section{Interpretacja wyników}
W jakim stopniu dane potwierdzają/obalają hipotezy.

\section{Porównanie z literaturą}
Zestawienie różnic (np. wzrost QUIC vs wcześniejsze badania).

\section{Zagrożenia dla ważności}
Internal: artefakty środowiska testowego; external: różnice regionów geograficznych; konstrukcja: wybór kategorii.

\section{Perspektywa praktyczna}
Wartość dla użytkowników, deweloperów, regulatorów.
\chapter{Podsumowanie i wnioski końcowe}

\section{Synteza wyników}

Celem pracy była empiryczna charakterystyka ruchu sieciowego wybranych aplikacji Android oraz ocena ryzyk prywatności i bezpieczeństwa na podstawie obserwowalnych metadanych (protokoły, hosty, porty, wolumeny) oraz analizy statycznej (trackery, uprawnienia). W toku badań:

\begin{itemize}
  \item zarejestrowano i zestandaryzowano sesje pracy aplikacji (Instagram, Facebook, Posnania, mObywatel, USOS, Spotify, Voronoi, TikTok),
  \item porównano protokoły transportowe (HTTP/1.1, HTTP/2, HTTP/3/QUIC, TLS), koncentrację ruchu na hostach oraz obecność telemetrii,
  \item skorelowano obserwacje z wynikami z Exodus Privacy (liczba trackerów i kategorie uprawnień).
\end{itemize}

\noindent Najważniejsze tendencje:
\begin{enumerate}
  \item \textbf{Aplikacje multimedialne/społecznościowe} (Instagram, Facebook, TikTok) wykazują wyraźną dominację \textbf{HTTP/3/QUIC} i dużą asymetrię RX/TX (konsumpcja treści).
  \item \textbf{Aplikacje instytucjonalne} (mObywatel, USOS) komunikują się konsekwentnie po \textbf{TLS/HTTPS} (brak QUIC, brak HTTP w klarze w zarejestrowanych oknach) i mają ograniczone odwołania do domen trzecich.
  \item \textbf{Aplikacje komercyjne} (Posnania, Spotify, Voronoi) opierają się głównie na \textbf{TLS/HTTPS} z typową telemetrią (np. Sentry, Crashlytics) i zapleczem w chmurze (AWS/Google).
\end{enumerate}

% \begin{table}[H]
% \centering
% \caption{Porównanie badanych aplikacji — zwięzły przegląd sesji i protokołów}
% \label{tab:apps-overview}
% \small
% \begin{tabular}{lrrrrll}
% \toprule
% \textbf{Aplikacja} & \multicolumn{1}{c}{\textbf{Flows}} & \multicolumn{1}{c}{\textbf{MiB}} & \multicolumn{1}{c}{\textbf{Okno}} & \textbf{Dominujący} & \textbf{HTTP/3} & \textbf{Uwagi (hosty/telemetria)} \\
% \midrule
% Instagram & 302 & 28{,}04 & $\sim$12{,}2 min & QUIC & Tak & graph/gateway.\,instagram, edge-mqtt (5222), silna asymetria \\
% Facebook  & 305 & 29{,}75 & $\sim$12{,}3 min & QUIC & Tak & graph/gateway, edge-mqtt; \textasciitilde98{,}5\% wolumenu po QUIC \\
% Posnania  & 21  & 0{,}08  & $\sim$2{,}9 min  & HTTPS & Nie & yourb.app; Sentry (ingest.\,sentry.io), Stripe \\
% mObywatel & 13  & 0{,}21  & $\sim$2{,}9 min  & HTTPS/TLS & Nie & api.\,mobywatel.gov.pl; sentry.\,puch.\,coi.\,gov.pl; jednorazowy HTTP/80 do Google \\
% USOS      & 200 & 2{,}88  & $\sim$4{,}3 min  & HTTPS & Nie & usosapps.\,put.\,poznan.pl, ppulse.\,put.\,poznan.pl; Crashlytics/Firebase \\
% Spotify   & 72  & 12{,}42 & $\sim$2{,}2 min  & HTTPS & Nie & spclient.\,wg.\,spotify.com; mDNS/SSDP (integracje urządzeń) \\
% Voronoi   & 465 & 9{,}61  & $\sim$5{,}6 min  & HTTPS & Nie & adresy AWS/Google (API/CDN); TLS/3000 śladowy \\
% TikTok    & --  & --      & --             & QUIC & Tak & domeny *tiktokv*.eu/\-tiktokcdn*, telemetria mcs/mon, silna asymetria \\
% \bottomrule
% \end{tabular}
% \vspace{0.5ex}
% \\ \footnotesize{Liczby \textit{Flows/MiB/Okno} odnoszą się do pojedynczych reprezentatywnych sesji zarejestrowanych w pracy; dla TikTok w części tabelarycznej nie zestawiano łącznych wartości, gdyż w rozdziale przedstawiono je głównie w ujęciu jakościowym.}
% \end{table}

\begin{table}[H]
\centering
\caption{Porównanie badanych aplikacji — zwięzły przegląd sesji i protokołów}
\label{tab:apps-overview}
\small
\setlength{\tabcolsep}{6pt}
\renewcommand{\arraystretch}{1.15}
\begin{tabularx}{\textwidth}{@{} l r r l l c X @{}}
\toprule
\textbf{Aplikacja} & \textbf{Flows} & \textbf{MiB} & \textbf{Okno} & \textbf{Dominujący} & \textbf{HTTP/3} & \textbf{Uwagi (hosty/telemetria)} \\
\midrule
Instagram & 302 & 28{,}04 & \(\sim\)12{,}2 min & QUIC       & Tak & graph/gateway.instagram, edge-mqtt (5222); silna asymetria RX/TX \\
Facebook  & 305 & 29{,}75 & \(\sim\)12{,}3 min & QUIC       & Tak & graph/gateway, edge-mqtt; \(\sim\)98{,}5\% wolumenu po QUIC \\
Posnania  & 21  & 0{,}08  & \(\sim\)2{,}9  min & HTTPS      & Nie & yourb.app; Sentry (ingest.sentry.io); Stripe \\
mObywatel & 13  & 0{,}21  & \(\sim\)2{,}9  min & HTTPS/TLS  & Nie & api.mobywatel.gov.pl; sentry.puch.coi.gov.pl; jednorazowy HTTP/80 do Google \\
USOS      & 200 & 2{,}88  & \(\sim\)4{,}3  min & HTTPS      & Nie & usosapps.put.poznan.pl, ppulse.put.poznan.pl; Crashlytics/Firebase \\
Spotify   & 72  & 12{,}42 & \(\sim\)2{,}2  min & HTTPS      & Nie & spclient.wg.spotify.com; mDNS/SSDP (integracje urządzeń) \\
Voronoi   & 465 & 9{,}61  & \(\sim\)5{,}6  min & HTTPS      & Nie & adresy AWS/Google (API/CDN); TLS/3000 śladowy \\
TikTok    & 649 & 125{,}30& \(\sim\)4{,}1  min & HTTPS (+ QUIC) & Tak & *tiktokv*.eu / tiktokcdn-eu.com; QUIC dla części strumieni; silna asymetria RX/TX \\
\bottomrule
\end{tabularx}
\vspace{0.25em}

{\footnotesize Dane pochodzą z pojedynczych sesji opisanych w rozdziale eksperymentów.}
\end{table}

\section{Wkład rozdziału „Pierwsze spostrzeżenia…”}

Przed badaniem właściwym wykonano serię prób i inwentaryzację, które bezpośrednio ukształtowały metodykę:

\paragraph{Inwentaryzacja aplikacji po resecie.}
Na dwóch urządzeniach po przywróceniu ustawień fabrycznych zidentyfikowano \textbf{53} domyślne aplikacje. Dla \textbf{44} (83{,}0\%) można było ręcznie \emph{zablokować} transfer w tle, dla \textbf{9} (17{,}0\%) — nie. Domyślnie \textbf{100\%} aplikacji miało \emph{aktywny} transfer w tle. Wniosek praktyczny: \textit{uprzednie} ograniczenie ruchu w tle (gdzie to możliwe) zmniejsza szum pomiarowy i poprawia rozdzielczość obserwacji scenariuszowych.

\begin{table}[H]
  \centering
  \caption{Ustawienia transferu danych w tle dla aplikacji domyślnych (po resecie)}
  \label{tab:bg-data-baseline}
  \begin{tabular}{lrr}
    \toprule
    \textbf{Kategoria} & \textbf{Liczba} & \textbf{Udział} \\
    \midrule
    Łączna liczba aplikacji            & 53 & 100\% \\
    Można zablokować dane w tle        & 44 & 83{,}0\% \\
    Nie można zablokować               & 9  & 17{,}0\% \\
    Domyślnie aktywny transfer w tle   & 53 & 100\% \\
    Domyślnie zablokowany transfer     & 0  & 0\% \\
    \bottomrule
  \end{tabular}
\end{table}

\paragraph{Przechwytywanie radiowe (Wireshark/monitor mode) — niepowodzenie.}
Sniffing ramek 802.11 z próbą dekrypcji ruchu aplikacji okazał się nieefektywny (złożone warunki radiowe, rotacja kluczy, brak pełnego materiału kryptograficznego). \textbf{Wniosek:} metoda ma wysoką cenę operacyjną i niską wartość poznawczą dla aplikacji mobilnych w warunkach produkcyjnych; lepiej \textit{kierować ruch przez własny węzeł} i rejestrować go na końcu.

\paragraph{MITM przez proxy (mitmproxy) — ograniczenia.}
Podejście z własnym AP i proxy ułatwia klasyfikację hostów i kontrolę DNS, ale w praktyce na przeszkodzie stają \textbf{pinning certyfikatów} i \textbf{HTTP/3/QUIC}, które omijają klasyczny „TLS-over-TCP” MITM. 

\paragraph{MITM + instrumentacja (Frida) — próba nieudana.}
Próba dynamicznego obejścia niestandardowych walidacji certyfikatów i pinningu zakończyła się niepowodzeniem (niestandardowe ścieżki weryfikacji, komponenty spoza prostych hooków). \textbf{Wniosek:} skuteczny MITM nowoczesnych aplikacji wymaga \emph{silnie customizowanej} instrumentacji lub środowisk buildowych — co wykracza poza zakres niniejszej pracy.

\paragraph{Rejestracja na urządzeniu (PCAP/PCAPNG) — sukces.}
Rejestracja bezpośrednio na telefonie (np. PCAPdroid) okazała się \textbf{najbardziej efektywna}, dając stabilny wgląd w metadane (protokoły/hosty/wolumeny) przy zachowaniu realnych warunków użytkowania, mimo braku semantyki HTTP przy QUIC/pinningu.

\section{Wnioski merytoryczne}

\begin{enumerate}
  \item \textbf{HTTP/3/QUIC jako standard mediów:} aplikacje z intensywną warstwą treści (Instagram, Facebook, TikTok) niemal całkowicie przeniosły strumienie na QUIC, co ogranicza możliwości klasycznego MITM i przesuwa analizę w stronę metadanych.
  \item \textbf{Transport TLS/HTTPS jako norma bezpieczeństwa:} w sesjach mObywatel i USOS nie odnotowano jawnego HTTP ani QUIC; ruch koncentrował się na hostach instytucjonalnych z niewielką telemetrią pomocniczą.
  \item \textbf{Telemetria w aplikacjach komercyjnych:} Posnania (Sentry, Stripe) i Spotify (ograniczona telemetria Google/Crashlytics) pokazują typowy profil integracji; Voronoi korzysta z chmury (AWS/Google) bez jawnego QUIC.
  \item \textbf{Ryzyka funkcjonalne niezależne od transportu:} w Posnanii zidentyfikowano możliwość \emph{dodania dowolnej tablicy rejestracyjnej} i wglądu w zdarzenia parkingowe — to obserwacja wymagająca dalszych badań on-device (nie przesądzana przez sam ruch sieciowy), ale istotna z perspektywy prywatności.
\end{enumerate}

\section{Ograniczenia badania}

\begin{itemize}
  \item wartości liczbowe dotyczą \emph{konkretnych wersji} aplikacji i \emph{pojedynczych okien czasowych}; zmiany wersji i regionu mogą wpływać na protokoły i hosty,
  \item QUIC i pinning ograniczają wgląd w semantykę HTTP — analiza dotyczy głównie metadanych,
  \item wyniki Exodus Privacy są statyczne (sygnatury obecności trackerów), a nie dowodem ich aktywności w danej sesji.
\end{itemize}

\section{Rekomendacje i kierunki dalszych prac}

\begin{enumerate}
  \item \textbf{Metodyka:} w badaniach terenowych preferować \emph{rejestrację on-device} + pasywną klasyfikację, a MITM stosować selektywnie (tam, gdzie brak QUIC i pinning pozwala na wgląd w HTTP).
  \item \textbf{Automatyzacja:} rozbudować pipeline o rozpoznawanie kategorii hostów (API/CDN/telemetria) i detekcję anomalii (np. nietypowe porty, nagły wzrost domen trzecich).
  \item \textbf{Prywatność:} dla aplikacji instytucjonalnych utrzymywać monitoring odwołań do domen trzecich i okresowo weryfikować politykę uprawnień.
  \item \textbf{Weryfikacja hipotez funkcyjnych:} przeprowadzić testy on-device (np. inspekcja bazy lokalnej, pinning, reverse engineering APK) dla przypadków takich jak rejestracje pojazdów w Posnanii.
\end{enumerate}

\section{Konkluzja}

Zastosowana metodyka pozwoliła na \textbf{spójne i powtarzalne} porównanie zachowań sieciowych różnych klas aplikacji. W praktyce użytkowej \textbf{transport jest w przeważającej mierze szyfrowany} (TLS/HTTPS lub QUIC), zaś realne ryzyka prywatności ujawniają się częściej w \emph{semantyce funkcji} i \emph{modelu danych} niż w „nagiej” warstwie protokołów. Badania potwierdzają kierunek: \textit{metadane są wystarczające do profilowania zachowań aplikacji}, ale \textit{ocena bezpieczeństwa} wymaga łączenia pasywnej obserwacji z testami dynamicznymi i analizą on-device.



% ======= BIBLIOGRAFIA =======
\backmatter
\printbibliography[heading=bibintoc,title={Bibliografia}]

% ======= DODATKI =======
\appendix
\chapter{Dodatkowe wykresy}
Dodatkowe wizualizacje (np. szczegółowe histogramy czasów odpowiedzi, rozkłady rozmiarów payloadów).
\chapter{Konfiguracja środowiska}
\section{Parametry mitmproxy}
\begin{verbatim}
mitmproxy --mode regular --set connection_strategy=lazy
\end{verbatim}

\section{Skrypt automatyzacji ADB (fragment)}
Przykład sekwencji wejść użytkownika generowanych sztucznie.

\section{Struktura katalogów danych}
Opis: raw/ processed/ reports/.

% Kolofon
\ppcolophon

\end{document}

% ======= BIBLIOGRAFIA =======
\backmatter
\printbibliography[heading=bibintoc,title={Bibliografia}]
